\documentclass[aspectratio=169]{beamer}
\usetheme{Madrid}
\useoutertheme{miniframes}
\usecolortheme{default}

\setbeamertemplate{headline}{
	\begin{beamercolorbox}[wd=\paperwidth,ht=3ex,dp=2.5ex]{section in head/foot}
		{\scriptsize \insertnavigation{\paperwidth}}
	\end{beamercolorbox}
}

\usepackage{amsmath,amssymb,bm}
\usepackage{mathtools}
\usepackage{booktabs}
\usepackage{listings}
\usepackage{xcolor}

\usepackage{tikz}
\usepackage{pgfplots}
\usepackage[most]{tcolorbox}
\pgfplotsset{compat=1.18}
\usetikzlibrary{arrows.meta,positioning,calc}


\title{Tutorial 13: EKF for Quadruped State Estimation}
\subtitle{Athletic Intelligence in Robotics (AIR) Course}
\author{Mihaela Popescu}
\date{06.05.2026}

\AtBeginSection[]{
	\begin{frame}[plain]
		\centering
		\tableofcontents[currentsection,sectionstyle=show/shaded,subsectionstyle=hide]
	\end{frame}
}

\begin{document}

\begin{frame}
	\titlepage
\end{frame}

\begin{frame}
	\frametitle{Outline}
	\tableofcontents
\end{frame}

\section{Introduction}
\begin{frame}
	\frametitle{Tutorial Goal}
	
	\begin{itemize}
		\item Implement an EKF for quadruped state estimation.
		\item Fuse:
		\begin{itemize}
			\item IMU measurements for prediction.
			\item Stance-foot contacts for correction.
		\end{itemize}
		\item Study the effect of:
		\begin{itemize}
			\item IMU noise
			\item IMU bias
			\item foot slip
			\item process and measurement noise tuning
			\item missing contacts / flight phases
		\end{itemize}
	\end{itemize}
	
	\begin{block}{Main idea}
		IMU integration predicts motion, while stance feet provide temporary world-frame anchors that reduce drift.
	\end{block}
	
\end{frame}

\section{EKF Description}
\begin{frame}
	\frametitle{State Representation in the EKF}
	
	\begin{itemize}
		\item The EKF maintains an estimate of the system state:
		\[
		\hat{x} =
		\left[
		R_{WB},\;
		v_W,\;
		p_W,\;
		p^{(1)}_{F,W},\;
		p^{(2)}_{F,W},\;
		\ldots
		\right]
		\]
		
		\item $R_{WB} \in \mathrm{SO}(3)$: rotation from body frame $\{B\}$ to world frame $\{W\}$.
		\item $v_W$: base velocity in world frame.
		\item $p_W$: base position in world frame.
		\item $p^{(i)}_{F,W}$: world-frame position of foot $i$ during stance.
		
		\vspace{0.2cm}
		
		\item The covariance is defined over an \textbf{error (perturbation) state}:
		\[
		\delta x =
		\left[
		\delta\theta,\;
		\delta v,\;
		\delta p,\;
		\delta p^{(1)}_F,\;
		\delta p^{(2)}_F,\;
		\ldots
		\right]
		\]
		
		\item $\delta\theta \in \mathbb{R}^3$ is a \textbf{minimal representation of orientation error} on $\mathrm{SO}(3)$.
	\end{itemize}
	
\end{frame}

\begin{frame}
	\frametitle{Contact-Dependent Foot States}
	
	\begin{itemize}
		\item Foot states are added only when a foot enters contact.
		\item When a foot is in stance, the no-slip assumption gives:
		\[
		p_{F,W} \approx \text{constant}
		\]
		
		\item When a new stance foot is detected, initialize its world position as:
		\[
		p_{F,W} = p_W + R_{WB}p_{F,B}
		\]
		
		\item Here:
		\begin{itemize}
			\item $p_{F,B}$ is the measured foot position in the body frame.
			\item $p_{F,W}$ is the stored foot anchor in the world frame.
		\end{itemize}
		
		\item When contact ends, the corresponding foot state is removed.
	\end{itemize}
	
	\begin{block}{Note}
		A swing foot should not be treated as fixed, because it moves relative to the world.
	\end{block}
	
\end{frame}


\begin{frame}
	\frametitle{Useful Operator: Skew Matrix}
	
	\begin{itemize}
		\item The skew matrix represents the cross product:
		\[
		[a]_\times b = a \times b
		\]
		
		\item For
		\[
		a =
		\begin{bmatrix}
			a_x \\ a_y \\ a_z
		\end{bmatrix},
		\]
		the skew matrix is:
		\[
		[a]_\times =
		\begin{bmatrix}
			0 & -a_z & a_y \\
			a_z & 0 & -a_x \\
			-a_y & a_x & 0
		\end{bmatrix}
		\]
		
		\item Small rotations are applied using the \textbf{matrix exponential}:
		\[
		R_{\text{new}} = R\exp([\delta\theta]_\times)
		\]
	\end{itemize}
	
\end{frame}


\begin{frame}
	\frametitle{Prediction Step: Orientation}
	
	\begin{itemize}
		\item The \textbf{gyroscope} measures angular velocity in the \textbf{body frame}:
		\[
		\omega_B =
		\begin{bmatrix}
			\omega_x & \omega_y & \omega_z
		\end{bmatrix}^T
		\]
		
		\item Over one time step $\Delta t$, compute:
		\[
		\Delta R = \exp([\omega_B]_\times \Delta t)
		\]
		
		\item Propagate orientation:
		\[
		R_{WB,k} = R_{WB,k-1}\Delta R
		\]
	\end{itemize}
	
\end{frame}


\begin{frame}
	\frametitle{Prediction Step: Acceleration, Velocity, Position}
	
	\begin{itemize}
		\item The \textbf{accelerometer} measures specific force in the \textbf{body frame}.
		\item Rotate it into the world frame and add gravity:
		\[
		a_W = R_{WB,k}a_B + g_W
		\]
		
		\item Propagate velocity:
		\[
		v_{W,k} = v_{W,k-1} + a_W\Delta t
		\]
		
		\item Propagate position:
		\[
		p_{W,k} =
		p_{W,k-1}
		+ v_{W,k-1}\Delta t
		+ \frac{1}{2}a_W\Delta t^2
		\]
		
	\end{itemize}
	
\end{frame}

\begin{frame}
	\frametitle{Prediction Step: Error-State Jacobian}
	
	\begin{itemize}
		\item Covariance propagation:
		\begin{equation*}
			P_k^- = F_k P_{k-1} F_k^T + Q
		\end{equation*}
		
		\item \textbf{Error-state vector}, where $\delta p_F$ are the stacked errors of all active feet:
		\begin{equation*}
			\delta x =
			\begin{bmatrix}
				\delta\theta & \delta v & \delta p & \delta p_F
			\end{bmatrix}^T
		\end{equation*}
		
	\item \textbf{Jacobian} from linearization ($F_{ab} = \frac{\partial a}{\partial b}$):
	\vspace{-0.2cm}
	\begin{columns}[T]
		\begin{column}{0.25\textwidth}
			\begin{equation*}
				F =
				\begin{bmatrix}
					F_{\theta\theta} & 0 & 0 & 0 \\
					F_{v\theta} & I & 0 & 0 \\
					F_{p\theta} & F_{pv} & I & 0 \\
					0 & 0 & 0 & I
				\end{bmatrix}
			\end{equation*}
		\end{column}
		
		\begin{column}{0.65\textwidth}
			\begin{equation*}
				\begin{aligned}
					F_{\theta\theta} &= \Delta R^T 
					&\qquad
					F_{v\theta} &= -R_{WB,k}[a_B]_\times \Delta t \\
					F_{p\theta} &= -\frac{1}{2}R_{WB,k}[a_B]_\times \Delta t^2 
					&\qquad
					F_{pv} &= I\,\Delta t
				\end{aligned}
			\end{equation*}
		\end{column}
	\end{columns}
		
		\item Process noise $Q$ captures model uncertainty:
		\begin{equation*}
			Q = \mathrm{diag}(Q_\theta,\; Q_v,\; Q_p)
		\end{equation*}
	\end{itemize}
	
\end{frame}



\begin{frame}
	\frametitle{Contact Update: Measurement Model}
	
	\begin{itemize}
		\item For a stance foot, predict the foot position from the current base estimate:
		\[
		\hat{p}_{F,W}
		=
		p_W + R_{WB}p_{F,B}
		\]
		
		\item Compare this with the stored foot anchor:
		\[
		r =
		\hat{p}_{F,W} - p_{F,W}
		\]
		
		\item If the foot is really fixed, this residual should be close to zero.
		
		\item The EKF update uses this residual to correct:
		\begin{itemize}
			\item base orientation,
			\item base position,
			\item foot anchor position.
		\end{itemize}
	\end{itemize}
	
	\begin{block}{Key assumption}
		The stance foot does not slip.
	\end{block}
	
\end{frame}


\begin{frame}
	\frametitle{Contact Update: Measurement Jacobian}
	
	\begin{itemize}
		\item The residual is:
		\[
		r =
		p_W + R_{WB}p_{F,B} - p_{F,W}
		\]
		
		\item Linearized with respect to the error state:
		\[
		H =
		\begin{bmatrix}
			H_\theta & H_v & H_p & H_{p_F}
		\end{bmatrix}
		\]
		
		\item Main blocks:
		\[
		H_\theta = -R_{WB}[p_{F,B}]_\times,\quad
		H_v = 0,\quad
		H_p = I,\quad
		H_{p_F} = -I
		\]
		
		\item \textit{Remark:} The orientation Jacobian $H_\theta$ can be unreliable in real world scenarios due to foot slip, imperfect contact detection, kinematic modeling errors, compliance and soft ground, frame misalignment, etc.
		
		\item For robustness, set $H_\theta = 0$ on real data. Position terms still constrain the base:
		\[
		H_p = I,\quad H_{p_F} = -I
		\]
	\end{itemize}
	
\end{frame}


\begin{frame}
	\frametitle{Contact Update: Kalman Correction}
	
	\begin{itemize}
		\item Innovation covariance:
		\[
		S = HP^-H^T + R
		\]
		
		\item Kalman gain:
		\[
		K = P^-H^TS^{-1}
		\]
		
		\item Error-state correction:
		\[
		\delta x = -Kr
		\]
		
		\item Inject the correction:
		\[
		R_{WB} \leftarrow R_{WB}\exp([\delta\theta]_\times)
		\]
		\[
		v_W \leftarrow v_W + \delta v
		\]
		\[
		p_W \leftarrow p_W + \delta p
		\]
		\[
		p_{F,W} \leftarrow p_{F,W} + \delta p_F
		\]
	\end{itemize}
	
\end{frame}


\begin{frame}
	\frametitle{Contact Update: Covariance}
	
	\begin{itemize}
		\item Use the Joseph-form covariance update:
		\[
		P =
		(I-KH)P^-(I-KH)^T
		+
		KRK^T
		\]
		
		\item This is more numerically stable than:
		\[
		P = (I-KH)P^-
		\]
		
		\item It helps preserve:
		\[
		P = P^T,
		\qquad
		P \succeq 0
		\]
		
		\item In code, symmetrize after the update:
		\[
		P \leftarrow \frac{1}{2}(P + P^T)
		\]
	\end{itemize}
	
\end{frame}


\begin{frame}
	\frametitle{What Can Go Wrong?}
	
	\begin{itemize}
		\item \textbf{IMU noise}: causes random drift after integration.
		
		\item \textbf{IMU bias}: causes systematic drift, often much worse than noise.
		
		\item \textbf{No contact}: the robot is temporarily dead-reckoning from IMU only.
		
		\item \textbf{Foot slip}: violates the fixed-contact assumption.
		
		\item \textbf{Wrong noise tuning}:
		\begin{itemize}
			\item small $Q$: filter trusts prediction too much
			\item large $R$: filter ignores foot updates
			\item small $R$: filter over-trusts possibly slipping contacts
		\end{itemize}
	\end{itemize}
	
\end{frame}

\section{Tutorial Tasks}
\begin{frame}
	\frametitle{Tutorial Workflow}
	
	\begin{enumerate}
		\item Run the provided simulation and inspect the data.
		\item Visualize the EKF state and active foot states.
		\item Implement IMU prediction.
		\item Test IMU-only estimation.
		\item Add IMU noise and bias.
		\item Implement foot-contact correction.
		\item Add foot slip and observe failure cases.
		\item Tune $Q$ and $R$.
		\item Visualize covariance ellipses.
		\item Apply the same estimator to real quadruped data.
	\end{enumerate}
	
\end{frame}


\begin{frame}
	\frametitle{Task 1--2: Data and State Inspection}
	
	\begin{itemize}
		\item Start with synthetic data. This provides ground truth for evaluation as follows:
		\[
		e_p(t) = p_{\text{est}}(t) - p_{\text{gt}}(t)
		\]
		\[
		\mathrm{RMSE}
		=
		\sqrt{
			\frac{1}{N}
			\sum_{k=1}^{N}
			\|e_p(k)\|^2
		}
		\]
		
		\item Inspect the initialized state:
		\[
		x =
		[R_{WB}, v_W, p_W]
		\]
		
		\item Then add stance-foot states using contact information:
		\[
		x =
		[R_{WB}, v_W, p_W, p^{(1)}_{F,W}, \ldots]
		\]
	\end{itemize}
	
\end{frame}


\begin{frame}
	\frametitle{Task 3: Implement IMU Propagation}
	
	\begin{enumerate}
		\item Propagate orientation:
		\[
		R_k = R_{k-1}\exp([\omega_k]_\times\Delta t)
		\]
		
		\item Compute world-frame acceleration:
		\[
		a_W = R_ka_B + g_W
		\]
		
		\item Integrate velocity:
		\[
		v_k = v_{k-1} + a_W\Delta t
		\]
		
		\item Integrate position:
		\[
		p_k =
		p_{k-1}
		+ v_{k-1}\Delta t
		+ \frac{1}{2}a_W\Delta t^2
		\]
		
		\item Build $F_k$ and propagate:
		\[
		P_k^- = F_kP_{k-1}F_k^T + Q
		\]
	\end{enumerate}
	
\end{frame}


\begin{frame}
	\frametitle{Task 4--5: IMU Noise and Bias}
	
	\begin{itemize}
		\item Add random IMU noise (zero-mean Gaussian noise $w_a, w_\omega$):
		\[
		a_m = a + w_a,
		\qquad
		\omega_m = \omega + w_\omega
		\]
		
		\item Add systematic IMU bias ($b_a, b_\omega$):
		\[
		a_m = a + b_a + w_a
		\]
		\[
		\omega_m = \omega + b_\omega + w_\omega
		\]
		
		\pause
		\textbf{Answers}
		\item Noise causes drift because acceleration and angular velocity are integrated.
		\item Bias causes larger long-term drift due to accumulation over time.
		
		\item Gyro bias is especially harmful:
		\[
		\omega_m = \omega + b_\omega
		\Rightarrow
		R_{WB} \text{ drifts over time}
		\Rightarrow
		a_W \text{ is rotated incorrectly}
		\Rightarrow
		p_W \text{ drifts}
		\]
	\end{itemize}
	
\end{frame}


\begin{frame}
	\frametitle{Task 6: Implement Foot-Contact Correction}
	
	\begin{enumerate}
		\item Compute predicted foot position:
		\[
		\hat{p}_{F,W} = p_W + R_{WB}p_{F,B}
		\]
		
		\item Compute residual:
		\[
		r = \hat{p}_{F,W} - p_{F,W}
		\]
		
		\item Build measurement Jacobian:
		\[
		H =
		\begin{bmatrix}
			-R_{WB}[p_{F,B}]_\times
			&
			0
			&
			I
			&
			-I
		\end{bmatrix}
		\]
		
		\item Compute:
		\[
		S = HP^-H^T + R,
		\qquad
		K = P^-H^TS^{-1}
		\]
		
		\item Correct state:
		\[
		\delta x = -Kr
		\]
		
		\item Update covariance using Joseph form.
	\end{enumerate}
	
\end{frame}


\begin{frame}
	\frametitle{Task 7: Foot Slip}
	
	\begin{itemize}
		\item Introduce foot slip in the simulation.
		
		\pause
		\textbf{Answer}
		\item The contact update assumes $p_{F,W} \approx \text{constant}$.		
		\item Foot slip violates this assumption:
		\[
		p_{F,W}(t + \Delta t) \neq p_{F,W}(t)
		\]
		
		\item Then the residual is no longer caused only by state-estimation error.
		\item The EKF may apply a wrong correction:
		\[
		\delta x = -Kr
		\]		
		\item Result:
		\begin{itemize}
			\item position error can increase,
			\item orientation can be pulled in the wrong direction,
			\item covariance can become overconfident,
			\item the estimate can become inconsistent.
		\end{itemize}
	\end{itemize}
	
\end{frame}


\begin{frame}
	\frametitle{Task 8: Tuning $Q$ and $R$}
	
	\begin{itemize}
		\item Process noise $Q$ controls trust in prediction:
		\[
		P^- = FP F^T + Q
		\]
		
		\item Measurement noise $R$ controls trust in contact updates:
		\[
		S = HP^-H^T + R
		\]
		
		\vspace{-0.4cm} \pause
		\textbf{Answer}
		\item If $Q$ is too small:
		\begin{itemize}
			\item the filter trusts the IMU prediction too much,
			\item covariance becomes too small,
			\item contact corrections have too little effect.
		\end{itemize}
		
		\item If $R$ is too large:
		\begin{itemize}
			\item the filter ignores foot contacts,
			\item behavior approaches IMU-only dead reckoning.
		\end{itemize}
		
		\item If $R$ is too small:
		\begin{itemize}
			\item the filter over-trusts contacts,
			\item foot slip can strongly corrupt the estimate.
		\end{itemize}
	\end{itemize}
	
\end{frame}


\begin{frame}
	\frametitle{Task 9: Covariance Ellipses}
	
	\begin{itemize}
		\item The position covariance in the $x$-$y$ plane is:
		\[
		P_{xy}
		=
		\begin{bmatrix}
			P_{xx} & P_{xy} \\
			P_{yx} & P_{yy}
		\end{bmatrix}
		\]
		
		\item Covariance ellipses visualize uncertainty.
		
		\pause
		\textbf{Answer}
		
		\item Ellipse size $\sim$ uncertainty (larger = higher, smaller = lower)		
		\item IMU-only prediction - $P$ grows; Contact updates - $P$ shrinks.		
		\item During flight phases, no foot anchor is available, so uncertainty grows.
	\end{itemize}
	
\end{frame}


\begin{frame}
	\frametitle{Task 10: Real Data}
	
	\begin{itemize}	
		
		\item Apply the same EKF for real world data and tune the $Q$ and $R$ noise parameters.
		
		\pause
		\textbf{Answer}
		\item Real data is harder than simulation because of:
		\begin{itemize}
			\item IMU bias and calibration errors,
			\item time synchronization errors,
			\item imperfect contact detection,
			\item compliance and impacts,
			\item kinematic errors,
			\item foot slip,
			\item uneven terrain.
		\end{itemize}
		
		\item For rough terrain, possible improvements include:
		\begin{itemize}
			\item contact probability estimation,
			\item slip detection,
			\item IMU bias estimation,
			\item adaptive measurement noise,
			\item visual or LiDAR odometry fusion.
		\end{itemize}
	\end{itemize}
	
\end{frame}

\section{Conclusion}
\begin{frame}
	\frametitle{Key Takeaways}
	
	\begin{itemize}
		\item EKF combines:
		\begin{itemize}
			\item IMU-based prediction (drift-prone),
			\item contact-based corrections (drift-reducing).
		\end{itemize}
		
		\item Stance feet act as temporary world-frame anchors.
		
		\item Performance depends strongly on:
		\begin{itemize}
			\item noise tuning ($Q$, $R$),
			\item contact quality (slip, detection),
			\item modeling assumptions.
		\end{itemize}
		
		\item Real-world data is more challenging than simulation.
		
		\item Error-state formulation enables efficient handling of orientation in $\mathrm{SO}(3)$.
	\end{itemize}
	
\end{frame}

\end{document}
